\chapter{CASSCF state data from the property file (menu 31)}
\label{ch:menu31}
\menulabel{31}

\section{Purpose}

Every ORCA job writes a \file{<base>.property.txt} with its computed
properties in a typed, sectioned format (the manual's property-file appendix
lists the schema; \code{orca\_2json <base> -property} converts the same file
to JSON).  For a CASSCF job the file carries the \emph{per-state} table
(\code{\$CAS\_SCF\_Energies}: total energy, multiplicity and irrep of every
state) and the \emph{absorption spectrum} section (state pairs, multiplicities
and excitation energies).  The menu reads those two sections and reports them;
this is also the landing point of the state-tracking survey (the
chapter's last section records the data situation for state identity along
a series).

\section{What you need}

The property file of a CASSCF run (written automatically; nothing to switch
on).  The state table needs \code{\$CAS\_SCF\_Energies}; the transition list
needs \code{CASSCF\_Absorption\_Spectrum}, which a run without the spectrum
request does not carry -- the menu then reports the state table alone.

\section{Walkthrough}

\lstinputlisting[style=fbkcode, firstline=54, lastline=74,
  title={The menu-31 example session (the N$_2$ SA-CASSCF property file: the
  per-state table, the transitions, the unnamed-column boundary and the
  availability survey)}]
  {examples/expected/m31_statedata.txt}

\section{What you get}

The per-state table (block, root, multiplicity, irrep, energy and the energy
relative to the lowest state in eV) and the absorption transitions (state
pair with irreps, multiplicities, $\Delta E$ in eV and cm$^{-1}$), plus a
report \file{<property>.states.fbk.md}.  The measured cross-checks: the
per-state energies reproduce the output's final \code{ROOT n: E=} lines to
their printed precision, and the transitions' first two columns satisfy the
eV/cm$^{-1}$ conversion (8065.544) as well as the state table's own energy
differences.  Columns the manual's schema leaves unnamed are printed in the
report under an explicit note, never interpreted.

\section{State tracking: what one output can and cannot say}

\begin{readbox}
\begin{itemize}
  \item \textbf{CI vectors are not persistable}: the CASSCF module keeps them
    in run-time temporaries, and the \file{.cis} file the JSON exporter looks
    for belongs to the CIS/STEOM modules (\code{orca\_2json} warns ``STEOM
    CIS file could not be found'' on a CASSCF job).  State \emph{identity}
    along a geometry series therefore cannot be read from one output;
  \item \textbf{what is printed}: the output's initial ``INITIAL CI STATE
    CHECK'' and the final ``CAS-SCF STATES FOR BLOCK'' blocks carry each
    root's dominant CSF occupations (two snapshots, before and after the
    orbital optimisation); this menu's property sections add the structured
    per-state energies and the transition list;
  \item \textbf{per-state observables for identity tracking}: single-root runs
    carry the per-state dipole (the menu-19 route), and the FIC-NEVPT2 sidecar
    carries per-root densities (the menu-23 chain) -- the tracker across a
    series built on those fingerprints is menu~\menu{49} (an ordered run
    sequence and a target root to the per-step root-matching table; its own
    chapter).  The state-averaged dipole in the property
    file is one x/y/z vector (\code{State -1}), not a per-root table: the
    menus-19/20 boundary stands, with that addition.
\end{itemize}
\end{readbox}
